\section{Methodology}

The Cascade Engine architecture consists of sequential triage layers, followed by a dynamic routing decision module.

\subsection{Intelligent Pre-Processing Layers}
Before a request reaches an LLM, it is processed by the following layers to ensure safety, efficiency, and context-awareness:
\begin{enumerate}
    \item \textbf{PrivacyFilter}: Utilizes Microsoft Presidio to identify and redact Personally Identifiable Information (PII), preventing sensitive data leaks to third-party APIs.
    \item \textbf{SemanticCache}: Uses FAISS and Sentence-Transformers to perform a vector similarity search against previously answered queries, short-circuiting the LLM entirely if a high-confidence match is found.
    \item \textbf{Gatekeeper}: A fine-tuned DistilBERT model that assesses the syntactic and semantic complexity of the query to preemptively route to the appropriate tier.
\end{enumerate}

\subsection{Routing Strategies}
If the request bypasses the semantic cache, it is forwarded to the Router. We implement several routing protocols:
\begin{itemize}
    \item \textbf{Frugal Routing}: Always attempts Tier 1 first, escalating to Tier 2 and Tier 3 only if the output confidence or validation metric is below a predefined threshold.
    \item \textbf{Learned Routing}: Employs a lightweight regression model to predict the required capability tier based on query embeddings and historical performance data.
\end{itemize}
